% ============================================================================
\section{Riemann--Hurwitz 公式与亏格}
\label{sec:riemann-hurwitz}
% ============================================================================

\textbf{为什么现在要问亏格？}

上一节我们把 $Y(\Gamma)$ 紧化成了紧黎曼面 $X(\Gamma)$。
对紧黎曼面来说，最重要的全局不变量之一就是亏格 $g$：
它衡量曲面有多少个"洞"，也控制全纯微分形式空间的维数。
换句话说，亏格不是一个可有可无的拓扑数字，
而是之后维数公式的几何心脏。

\textbf{核心想法是拿 $X(\Gamma)$ 去和最熟悉的曲面比较。}
最熟悉的紧黎曼面是
$X(1)\cong\mathbb{P}^1(\C)$，它的亏格是 $0$。
而包含关系 $\Gamma\subset \SL_2(\Z)$ 给出自然映射
\[
  f\colon X(\Gamma)\longrightarrow X(1).
\]
这个映射在绝大多数点看起来像普通覆盖，
只有在椭圆点和尖点上会发生分歧。
Riemann--Hurwitz 公式正是把"覆盖的度数"和"分歧有多严重"翻译成亏格的公式。

\textbf{得到的回报非常大。}
一旦把分歧项数清楚，
$g\bigl(X(\Gamma)\bigr)$ 就能显式写出来；
特别对 $\Gamma_0(N)$，这给出一条可以真正计算的 genus formula，
而这正是后面讨论 $\dim \Mk_k(\Gamma)$ 与 $\dim \Sk_k(\Gamma)$ 的起点。

\begin{theorem}[Riemann--Hurwitz 公式，回顾]
\label{thm:riemann-hurwitz-recall}
设 $f\colon X\to Y$ 是紧黎曼面之间度为 $d$ 的非常值全纯映射。
记 $e_P$ 为 $f$ 在 $P\in X$ 处的分歧指数，则
\[
  2g(X)-2=d\bigl(2g(Y)-2\bigr)+\sum_{P\in X}(e_P-1).
\]
\end{theorem}

\begin{remark}
严格证明已在\cref{thm:riemann-hurwitz} 给出。
本节的重点不是再证明一次，而是把它真正用到模曲线上。
\end{remark}

% figures/ch02-branched-cover.tex
% 模曲线到 X(1)\cong P^1 的分歧覆盖示意图
\begin{figure}[ht]
\centering
\begin{tikzpicture}[>=Stealth, font=\small, scale=1.0]
  % top surface
  \draw[thick, fill=blue!8] (-1.8,2.2) .. controls (-0.8,3.2) and (1.2,3.3) .. (2.1,2.5)
    .. controls (3.0,1.8) and (2.6,0.7) .. (1.5,0.3)
    .. controls (0.2,-0.2) and (-1.3,0.2) .. (-1.8,1.2) -- cycle;
  \node at (0.25,1.8) {$X(\Gamma)$};
  \fill[red!80!black] (-0.7,2.15) circle (2pt);
  \fill[red!80!black] (1.35,2.45) circle (2pt);
  \fill[green!60!black] (0.1,0.75) circle (2pt);
  \node[above left, red!80!black] at (-0.7,2.15) {上方的 $i$};
  \node[above right, red!80!black] at (1.35,2.45) {上方的 $\rho$};
  \node[below, green!60!black] at (0.1,0.75) {上方的尖点};

  % arrows
  \draw[->, very thick] (-0.3,0.15) -- (-0.3,-1.0);
  \draw[->, very thick] (0.7,0.15) -- (0.7,-1.0);
  \node[right] at (0.7,-0.45) {$f\colon X(\Gamma)\to X(1)$};

  % sphere / projective line
  \draw[thick, fill=orange!10] (-2.0,-2.7) .. controls (-0.9,-3.5) and (1.4,-3.5) .. (2.2,-2.6)
    .. controls (3.0,-1.7) and (2.2,-0.8) .. (0.4,-0.7)
    .. controls (-1.2,-0.7) and (-2.4,-1.5) .. (-2.0,-2.7) -- cycle;
  \node at (0.2,-1.95) {$X(1)\cong \mathbb{P}^1$};
  \fill[red!80!black] (-0.95,-1.25) circle (2.2pt);
  \fill[red!80!black] (1.2,-1.2) circle (2.2pt);
  \fill[green!60!black] (0.15,-2.7) circle (2.2pt);
  \node[above, red!80!black] at (-0.95,-1.25) {$i$};
  \node[above, red!80!black] at (1.2,-1.2) {$\rho$};
  \node[below, green!60!black] at (0.15,-2.7) {$\infty$};

  \node[align=center] at (4.4,-1.9)
    {Riemann--Hurwitz 要做的事：\\
     度数 $\mu$ 给出主项，\\
     分歧点 $i,\rho,\infty$ 给出修正项。};
\end{tikzpicture}
\caption{紧化模曲线到 $X(1)\cong\mathbb{P}^1$ 的覆盖示意图。分歧现象集中在椭圆点和尖点，它们正是亏格公式中的修正项来源。}
\label{fig:branched-cover-modular-curve}
\end{figure}


\begin{definition}[模曲线亏格公式中的记号]
\label{def:genus-formula-notation}
设 $\Gamma$ 为有限指数子群，记 $\bar\Gamma$ 为它在 $\PSL_2(\Z)$ 中的像，并设
\[
  \mu=[\PSL_2(\Z):\bar\Gamma].
\]
若 $-I\in\Gamma$（例如 $\Gamma=\Gamma_0(N)$），则 $\mu=[\SL_2(\Z):\Gamma]$。
再记
\begin{itemize}
  \item $\nu_2$：$Y(\Gamma)$ 上阶 $2$ 椭圆点的个数；
  \item $\nu_3$：$Y(\Gamma)$ 上阶 $3$ 椭圆点的个数；
  \item $\nu_\infty$：尖点个数。
\end{itemize}
\end{definition}

\begin{theorem}[模曲线的 genus formula]
\label{thm:genus-formula-modular-curves}
设 $\Gamma$ 为有限指数子群并且 $-I\in\Gamma$。
则紧化模曲线 $X(\Gamma)$ 的亏格满足
\[
  g\bigl(X(\Gamma)\bigr)
  =1+\frac{\mu}{12}-\frac{\nu_2}{4}-\frac{\nu_3}{3}-\frac{\nu_\infty}{2}.
\]
特别地，对 $\Gamma_0(N)$，这里的 $\mu$ 就是
\[
  \mu=[\SL_2(\Z):\Gamma_0(N)]=N\prod_{p\mid N}\left(1+\frac1p\right).
\]
\end{theorem}

\begin{proof}
把 Riemann--Hurwitz 应用于
\[
  f\colon X(\Gamma)\longrightarrow X(1)\cong\mathbb{P}^1.
\]
由于 $g\bigl(X(1)\bigr)=0$，公式变成
\[
  2g\bigl(X(\Gamma)\bigr)-2=-2\mu+\sum_{P\in X(\Gamma)}(e_P-1).
\]
所以关键是计算总分歧贡献。

\textbf{(1) 在 $i$ 上方的贡献。}
$X(1)$ 在 $i$ 处的局部稳定子阶为 $2$。
因此 $f^{-1}(i)$ 中的点分两类：
若该点在 $X(\Gamma)$ 上是阶 $2$ 椭圆点，则局部度数为 $1$；
否则局部度数为 $2$，贡献 $e_P-1=1$。
故 $i$ 上方共有
\[
  \nu_2+\frac{\mu-\nu_2}{2}=\frac{\mu+\nu_2}{2}
\]
个点，总分歧贡献为
\[
  \frac{\mu-\nu_2}{2}.
\]

\textbf{(2) 在 $\rho=e^{2\pi i/3}$ 上方的贡献。}
同理，$\rho$ 的稳定子阶为 $3$。
阶 $3$ 椭圆点的局部度数为 $1$；其余点局部度数为 $3$，每点贡献 $2$。
所以总贡献为
\[
  2\cdot\frac{\mu-\nu_3}{3}=\frac{2(\mu-\nu_3)}{3}.
\]

\textbf{(3) 在尖点上方的贡献。}
设各尖点宽度为 $h_1,\dots,h_{\nu_\infty}$。
由于 $f$ 在尖点的局部坐标中形如 $q\mapsto q^{h_j}$，
第 $j$ 个尖点的分歧指数就是 $h_j$，贡献 $h_j-1$。
而所有宽度之和等于覆盖度数：
\[
  h_1+\cdots+h_{\nu_\infty}=\mu.
\]
因此尖点总贡献为
\[
  \sum_{j=1}^{\nu_\infty}(h_j-1)=\mu-\nu_\infty.
\]

把三部分相加：
\[
  \sum_P(e_P-1)=\frac{\mu-\nu_2}{2}+\frac{2(\mu-\nu_3)}{3}+\mu-\nu_\infty.
\]
代回 Riemann--Hurwitz 并整理，得到
\[
  g\bigl(X(\Gamma)\bigr)
  =1+\frac{\mu}{12}-\frac{\nu_2}{4}-\frac{\nu_3}{3}-\frac{\nu_\infty}{2}.
\qedhere
\]
\end{proof}

\begin{example}[三个最重要的例子]
\label{ex:genus-small-levels}
\begin{enumerate}
  \item $X_0(1)=X(1)$ 的亏格是 $0$，因为它就是 $\mathbb{P}^1$。
  \item $X_0(11)$：
    \[
      \mu=11\left(1+\frac1{11}\right)=12,
      \qquad \nu_2=0,
      \qquad \nu_3=0,
      \qquad \nu_\infty=2.
    \]
    所以
    \[
      g\bigl(X_0(11)\bigr)=1+\frac{12}{12}-0-0-\frac22=1.
    \]
    这意味着 $X_0(11)$ 是一条亏格 $1$ 的曲线——也就是一条椭圆曲线。
  \item $X_0(23)$：
    \[
      \mu=23\left(1+\frac1{23}\right)=24,
      \qquad \nu_2=0,
      \qquad \nu_3=0,
      \qquad \nu_\infty=2.
    \]
    因而
    \[
      g\bigl(X_0(23)\bigr)=1+\frac{24}{12}-1=2.
    \]
\end{enumerate}
\end{example}

\begin{remark}
请特别留意这个公式里每一项的几何来源：
$\mu/12$ 来自覆盖的主项，
$\nu_2/4$ 与 $\nu_3/3$ 来自椭圆点的对称性修正，
$\nu_\infty/2$ 来自尖点。
这正是模曲线把群作用、复分析和拓扑缝在一起的地方。
\end{remark}

\textbf{覆盖理论视角。}
把模曲线之间的自然态射看成“忘掉一部分级别结构”，
Riemann--Hurwitz 的分歧项就有了非常直观的模解释。

\begin{proposition}[遗忘级别结构的分歧]
\label{prop:forget-level-ramification}
若 $M\mid N$，则忘掉多余级别结构给出有限态射
\[
  \pi_{N,M}:X_0(N)\longrightarrow X_0(M),
\]
其度数为 $[\Gamma_0(M):\Gamma_0(N)]$。
在非椭圆点、非尖点处，$\pi_{N,M}$ 局部上是双全纯的；
全部分歧都来自两类点：
\begin{enumerate}
  \item 椭圆点：源点对应的椭圆曲线带有额外自同构，
        所以不同的级别结构在遗忘后会被识别；
  \item 尖点：源尖点与目标尖点的宽度不同。
        若该点在 $X_0(N)$、$X_0(M)$ 上的宽度分别为 $h_N,h_M$，
        则局部参数满足
        \[
          q_M=q_N^{h_N/h_M},
        \]
        因而分歧指数恰为 $h_N/h_M$。
\end{enumerate}
\end{proposition}

\begin{proof}[说明]
忘掉一部分循环子群显然给出一个有限覆盖：
对一般点而言，纤维只是同一条椭圆曲线上若干种补充级别结构的有限集合。
当椭圆曲线没有额外自同构时，这些选择彼此分开，故局部上没有分歧。
若底点恰是 $j=1728$ 或 $j=0$ 那样的椭圆点，则自同构群更大，
若干个级别结构会在遗忘后合并，于是出现椭圆点分歧。
对尖点则使用上一节的 $q$-坐标：
遗忘级别结构只会改变尖点宽度，故局部模型必为幂映射 $q\mapsto q^e$，
其指数正是宽度之比。
\end{proof}

\begin{corollary}[映到 $X(1)$ 时的三类分歧]
\label{cor:ramification-to-x1}
对自然映射
\[
  f:X_0(N)\longrightarrow X(1),
\]
全部分歧都落在 $j=1728,0,\infty$ 这三个特殊值上。
若把三处的总分歧贡献分别记为
\[
  e_2=\frac{\mu-\nu_2}{2},
  \qquad
  e_3=\frac{2(\mu-\nu_3)}{3},
  \qquad
  e_\infty=\mu-\nu_\infty,
\]
则 Riemann--Hurwitz 正好写成
\[
  2g\bigl(X_0(N)\bigr)-2=-2\mu+e_2+e_3+e_\infty.
\]
换言之，genus formula 里的三项修正分别来自阶 $2$ 椭圆点、阶 $3$ 椭圆点和尖点。
\end{corollary}

\begin{proof}[说明]
$X(1)$ 上除了 $i$、$\rho$ 与尖点 $\infty$ 之外，稳定子群都是平凡的，
所以映射在其余点上不可能分歧。
前面证明 genus formula 时已经逐项算过三类点的总贡献：
$i$ 上方给出 $(\mu-\nu_2)/2$，
$\rho$ 上方给出 $2(\mu-\nu_3)/3$，
尖点上方给出 $\mu-\nu_\infty$。
把它们重新记作 $e_2,e_3,e_\infty$，就得到上式。
\end{proof}

\textbf{各项的几何意义。}
从 CSS 中 Stevens 的视角看，
\begin{itemize}
  \item $\mu/12$ 是把 $X_0(N)$ 当成次数 $\mu$ 的覆盖后得到的主项；
  \item $\nu_2/4$ 与 $\nu_3/3$ 是对椭圆点对称性的“返还项”：
        稳定子越大，真正需要计入的分歧就越少；
  \item $\nu_\infty/2$ 来自无穷远处的退化纤维，也就是尖点的宽度分布。
\end{itemize}
因此亏格公式并不是孤立的代数恒等式，
而是把覆盖次数、局部对称性和尖点退化三种几何现象压缩成一行数字。
